%=============================================================================
% PATENT 15 — COMPLETE ENGINEERING DESIGN DOCUMENT
% "Laboratory Scalar Density Field Generator: Deployment-Ready
%  Engineering Design Specification"
%
% Supporting: US Provisional Patent 63/868,073 (Aug 21, 2025)
%=============================================================================
\documentclass[11pt,a4paper]{article}

\usepackage[margin=1in]{geometry}
\usepackage{amsmath,amssymb}
\usepackage{siunitx}
\usepackage{graphicx}
\usepackage{booktabs}
\usepackage{longtable}
\usepackage{hyperref}
\usepackage{xcolor}
\usepackage{enumitem}
\usepackage{tcolorbox}
\usepackage{fancyhdr}
\usepackage{titlesec}

\pagestyle{fancy}
\fancyhead[L]{\small Patent 15 --- Engineering Design}
\fancyhead[R]{\small CONFIDENTIAL}
\fancyfoot[C]{\thepage}

\title{\textbf{Laboratory Scalar Density Field Generator}\\[6pt]
\large Deployment-Ready Engineering Design Specification\\[6pt]
\normalsize Supporting US Provisional Patent 63/868,073\\
Document Revision 1.0}

\author{DFD Research Programme --- Patent 15 Technical Team}
\date{March 2026}

\begin{document}
\maketitle
\tableofcontents
\newpage

%=============================================================================
\section{Executive Summary}
%=============================================================================

This document provides the complete, deployment-ready engineering
design for a Laboratory Scalar Density Field Generator (LSDFG),
the enabling apparatus for all Density Field Dynamics (DFD)
technology patents.  The LSDFG generates controlled perturbations
of the scalar field $\psi$ through electromagnetic--$\psi$
back-reaction coupling, measures the generated field through a
multi-modal sensor suite, and stabilizes the field via real-time
feedback control.

\paragraph{Key specifications:}
\begin{itemize}[nosep]
\item \textbf{Detection sensitivity}: $\Delta\psi \sim 10^{-19}$
(Phase 3, $10^4\;\si{s}$ integration)
\item \textbf{$\lambda$ parameter reach}:
$|\lambda - 1| \sim 10^{-14}$ (Phase 3)
\item \textbf{Footprint}: $5\;\si{m} \times 5\;\si{m} \times 6\;\si{m}$
(including cryogenics)
\item \textbf{Power}: $\sim 200\;\si{kW}$ (dominated by cryogenics)
\item \textbf{Cooldown time}: $\sim 48\;\si{h}$ to operational
temperature
\item \textbf{Campaign duration}: $\sim 2$ weeks for Phase 1
broadband scan
\end{itemize}

%=============================================================================
\section{System Requirements}
%=============================================================================

\subsection{Functional Requirements}

\begin{longtable}{p{0.08\textwidth}p{0.60\textwidth}p{0.22\textwidth}}
\toprule
\textbf{ID} & \textbf{Requirement} & \textbf{Verification} \\
\midrule
\endhead
FR-01 & Generate EM-driven $\psi$ perturbations via TE+TM
superposition in SRF cavities & Analysis + Test \\
FR-02 & Detect $\psi$ perturbations at $\Delta\psi < 10^{-15}$
in $10^3\;\si{s}$ & Test \\
FR-03 & Measure $\psi$ independently in optical, atomic, and
inertial domains & Test \\
FR-04 & Stabilize generated $\psi$ field to $\pm 10\%$ of
setpoint & Test \\
FR-05 & Scan EM drive frequency over $0.1$--$10\;\si{GHz}$
range & Test \\
FR-06 & Enforce hardware safety limits on $|\psi|$,
$|\nabla\psi|$, and $d|\psi|/dt$ & Inspection + Test \\
FR-07 & Maintain vacuum $< 10^{-6}\;\si{Pa}$ during
operation & Test \\
FR-08 & Maintain thermal stability $\pm 1\;\si{mK}$ at
sensor locations & Test \\
FR-09 & Provide phase-coherent timing across all subsystems
at $< 100\;\si{ps}$ & Test \\
FR-10 & Record all sensor data with $< 100\;\si{ps}$
timestamps continuously & Inspection \\
\bottomrule
\end{longtable}

\subsection{Performance Requirements}

\begin{longtable}{p{0.08\textwidth}p{0.45\textwidth}p{0.18\textwidth}p{0.19\textwidth}}
\toprule
\textbf{ID} & \textbf{Parameter} & \textbf{Phase 1} & \textbf{Phase 3} \\
\midrule
\endhead
PR-01 & Stored EM energy & $9\;\si{kJ}$ & $1\;\si{MJ}$ \\
PR-02 & Cavity $Q_0$ & $> 10^{10}$ & $> 3 \times 10^{10}$ \\
PR-03 & Modulation depth $m$ & $0.01$ & $0.10$ \\
PR-04 & $\psi$ detection floor & $10^{-17}$ & $10^{-19}$ \\
PR-05 & Accelerometer noise & $10^{-11}\;\si{m/s^2/\sqrt{Hz}}$ &
$10^{-13}\;\si{m/s^2/\sqrt{Hz}}$ \\
PR-06 & Clock instability & $5 \times 10^{-17}/\sqrt{\tau}$ &
$2 \times 10^{-18}/\sqrt{\tau}$ \\
PR-07 & MZI phase noise & $10^{-9}\;\si{rad/\sqrt{Hz}}$ &
$10^{-11}\;\si{rad/\sqrt{Hz}}$ \\
\bottomrule
\end{longtable}

\subsection{Environmental Requirements}

\begin{longtable}{p{0.08\textwidth}p{0.52\textwidth}p{0.30\textwidth}}
\toprule
\textbf{ID} & \textbf{Requirement} & \textbf{Specification} \\
\midrule
\endhead
ER-01 & Ambient temperature & $20 \pm 2\;^\circ$C \\
ER-02 & Relative humidity & $30$--$60\%$ \\
ER-03 & Ambient vibration (floor) & $< 10^{-7}\;\si{m/s^2/\sqrt{Hz}}$
at $1\;\si{Hz}$ \\
ER-04 & Ambient magnetic field & $< 50\;\si{\mu T}$ (Earth's field) \\
ER-05 & Power supply & $3$-phase $400\;\si{V}$, $500\;\si{A}$ \\
ER-06 & Cooling water & $20\;\si{L/min}$ at $15\;^\circ$C \\
ER-07 & Liquid He supply & $500\;\si{L/day}$ (or closed-cycle) \\
ER-08 & Floor loading & $> 5000\;\si{kg/m^2}$ \\
ER-09 & Ceiling height & $> 6\;\si{m}$ \\
\bottomrule
\end{longtable}

%=============================================================================
\section{System Architecture}
%=============================================================================

\subsection{Block Diagram}

The LSDFG consists of five primary subsystems connected as follows:

\begin{tcolorbox}[colback=blue!3, colframe=blue!50!black,
title=System Architecture Block Diagram]
\begin{verbatim}
+-----------------------------------------------------------+
|                    CONTROL & DAQ (Sec 7)                  |
|  +----------+  +----------+  +----------+  +----------+  |
|  | Kalman   |  | Feedback |  | Safety   |  | Data     |  |
|  | Filter   |  | Control  |  | Interlock|  | Recorder |  |
|  +----+-----+  +----+-----+  +----+-----+  +----+-----+  |
|       |             |             |             |          |
+-------+-------------+-------------+-------------+---------+
        |             |             |             |
   +----v----+   +----v----+   +----v----+   +---v----+
   | SENSOR  |   | SOURCE  |   | VACUUM  |   | TIMING |
   | SUITE   |   | ARRAY   |   | & ENV.  |   | DIST.  |
   | (Sec 5) |   | (Sec 4) |   | (Sec 3) |   | (Sec 6)|
   +---------+   +---------+   +---------+   +--------+
\end{verbatim}
\end{tcolorbox}

\subsection{Phased Implementation}

\begin{longtable}{p{0.12\textwidth}p{0.25\textwidth}p{0.25\textwidth}p{0.28\textwidth}}
\toprule
\textbf{Phase} & \textbf{Sources} & \textbf{Sensors} & \textbf{Reach} \\
\midrule
\endhead
\textbf{Phase 1} (Yr 1) & 4$\times$ SRF cavities, CW, $U_0 = 9\;\si{kJ}$
& MZI + 2$\times$ FP cavities + 3$\times$ gravimeters
& $|\lambda{-}1| \sim 10^{-9}$ \\[6pt]
\textbf{Phase 2} (Yr 2) & Phase 1 + optimized array, rotating mass
& Phase 1 + Sr lattice clock + Rb atom interferometer
& $|\lambda{-}1| \sim 10^{-11}$ \\[6pt]
\textbf{Phase 3} (Yr 3+) & Pulsed SRF to $1\;\si{MJ}$, shrunk tube area
& Phase 2 + Yb clock + Th-229 nuclear clock
& $|\lambda{-}1| \sim 10^{-14}$ \\
\bottomrule
\end{longtable}

%=============================================================================
\section{Vacuum Chamber and Environmental Control}
%=============================================================================

\subsection{Primary Vacuum Vessel}

\subsubsection{Geometry and Materials}

\begin{longtable}{p{0.35\textwidth}p{0.55\textwidth}}
\toprule
\textbf{Parameter} & \textbf{Specification} \\
\midrule
\endhead
Overall geometry & Vertical cylinder, domed ends \\
Inner diameter & $1200\;\si{mm}$ \\
Cylindrical height & $3000\;\si{mm}$ \\
Total internal volume & $3.39\;\si{m^3}$ \\
Wall material & 316L stainless steel (low magnetic permeability) \\
Wall thickness & $12\;\si{mm}$ (ASME VIII Div.~1 rated) \\
End caps & Torispherical, $12\;\si{mm}$ \\
Surface finish (internal) & Electropolished, $R_a < 0.4\;\si{\mu m}$ \\
Surface finish (external) & 2B mill finish \\
Leak rate & $< 10^{-9}\;\si{Pa\cdot m^3/s}$ (He) \\
Design pressure & Vacuum to $0.2\;\si{MPa}$ (gauge) \\
Bakeout temperature & $150\;^\circ$C (full vessel) \\
Mass (empty) & $\sim 2500\;\si{kg}$ \\
\bottomrule
\end{longtable}

\subsubsection{Port Configuration}

\begin{longtable}{p{0.15\textwidth}p{0.15\textwidth}p{0.15\textwidth}p{0.45\textwidth}}
\toprule
\textbf{Port} & \textbf{Size} & \textbf{Qty} & \textbf{Purpose} \\
\midrule
\endhead
Optical & CF-100 & 12 & Laser access at $30^\circ$ azimuthal spacing \\
RF waveguide & WR-90 & 4 & SRF cavity feeds \\
RF coaxial & SMA & 8 & Monitoring, calibration, auxiliary feeds \\
Fiber & FC/APC & 6 & PM fiber for clock/interferometer links \\
Electrical & D-sub 50 & 4 & Sensor signals, thermocouples, heaters \\
Electrical & BNC & 24 & Fast analog signals \\
Pumping & DN-160 CF & 2 & Turbo pump, ion pump \\
Cryogenic & Custom & 4 & LHe transfer lines \\
Top access & DN-400 CF & 1 & Equipment installation \\
Bottom access & DN-250 CF & 1 & Feedthrough plate \\
\bottomrule
\end{longtable}

\subsection{Vacuum System}

\subsubsection{Pumping Chain}

\begin{enumerate}[nosep]
\item \textbf{Roughing}: Dry scroll pump, $10\;\si{m^3/h}$,
base $1\;\si{Pa}$.
\item \textbf{High vacuum}: Turbomolecular pump, $300\;\si{L/s}$
(N$_2$), magnetically levitated bearings, vibration
$< 10^{-3}\;\si{m/s^2}$.
\item \textbf{Ultra-high vacuum}: Sputter-ion pump, $60\;\si{L/s}$,
no moving parts, for operational phase.
\item \textbf{Pressure monitoring}: Bayard-Alpert ionization gauge
($10^{-8}$--$10^{-2}\;\si{Pa}$), Pirani gauge
($10^{-1}$--$10^5\;\si{Pa}$), residual gas analyzer
(1--200 AMU).
\end{enumerate}

\paragraph{Pumpdown schedule:}
Roughing to $1\;\si{Pa}$: $\sim 2\;\si{h}$.
Turbo to $10^{-5}\;\si{Pa}$: $\sim 12\;\si{h}$.
Bakeout + ion pump to $< 10^{-7}\;\si{Pa}$: $\sim 48\;\si{h}$.

\subsection{Thermal Control}

\subsubsection{Three-Stage Temperature Stabilization}

\begin{longtable}{p{0.12\textwidth}p{0.20\textwidth}p{0.20\textwidth}p{0.38\textwidth}}
\toprule
\textbf{Stage} & \textbf{Stability} & \textbf{Method} & \textbf{Details} \\
\midrule
\endhead
Outer & $\pm 0.1\;\si{K}$ & Water jacket & Cu tubing brazed to vessel exterior,
$5\;\si{L/min}$ chilled water loop \\
Middle & $\pm 1\;\si{mK}$ & Active Peltier & 16$\times$ TEC modules on Al shield,
PID-controlled to $10\;\si{\mu K}$ readout \\
Inner & $\pm 0.1\;\si{mK}$ & Passive radiation & 3 polished Al shields,
emissivity $< 0.05$, passive equilibration \\
\bottomrule
\end{longtable}

\paragraph{Temperature sensors:} 24$\times$ calibrated Pt-100 RTDs
(4-wire, $0.1\;\si{mK}$ resolution) distributed across all thermal
stages and critical sensor locations.  2$\times$ sapphire-diode
sensors for cryogenic stages.

\subsection{Magnetic Shielding}

Three concentric cylindrical mu-metal shields:
\begin{itemize}[nosep]
\item Layer 1 (outer): $1.5\;\si{mm}$ mu-metal, $\varnothing 1500\;\si{mm}$
\item Layer 2 (middle): $1.0\;\si{mm}$ mu-metal, $\varnothing 1400\;\si{mm}$
\item Layer 3 (inner): $0.8\;\si{mm}$ mu-metal, $\varnothing 1300\;\si{mm}$
\item Combined shielding factor: $> 10^5$ at DC, $> 10^6$ at $50\;\si{Hz}$
\item Residual field at center: $< 1\;\si{nT}$
\item Degaussing coils on each layer for periodic demagnetization
\end{itemize}

\subsection{Vibration Isolation}

\subsubsection{Stage 1: Inertial Block}

\begin{itemize}[nosep]
\item Reinforced concrete block: $2.5\;\si{m} \times 2.5\;\si{m}
\times 1.0\;\si{m}$, mass $\sim 15{,}000\;\si{kg}$
\item Supported on 4$\times$ pneumatic isolators (air springs)
\item Natural frequency: $f_0 = 0.5\;\si{Hz}$ (vertical),
$0.3\;\si{Hz}$ (horizontal)
\item Passive isolation: $-40\;\si{dB/decade}$ above $f_0$
\end{itemize}

\subsubsection{Stage 2: Active Platform}

\begin{itemize}[nosep]
\item Active piezo-servo platform beneath vacuum vessel
\item 6-DOF feedback using integrated geophones
\item Bandwidth: $0.5$--$200\;\si{Hz}$
\item Residual vibration: $< 10^{-9}\;\si{m/s^2/\sqrt{Hz}}$
above $1\;\si{Hz}$
\end{itemize}

\subsubsection{Stage 3: Internal Suspensions}

\begin{itemize}[nosep]
\item Critical sensor assemblies (FP cavities, clocks) on
monolithic fused-silica suspensions
\item Pendulum frequency: $f_0 = 0.1\;\si{Hz}$
\item Quality factor: $Q > 10^6$ (fused silica)
\item Effective isolation at $1\;\si{Hz}$: $> 10^4$
\end{itemize}

%=============================================================================
\section{Source Array Design}
%=============================================================================

\subsection{SRF Cavity Subsystem}

\subsubsection{Individual Cavity Design}

\begin{longtable}{p{0.40\textwidth}p{0.50\textwidth}}
\toprule
\textbf{Parameter} & \textbf{Value} \\
\midrule
\endhead
Cavity type & 9-cell TESLA-type elliptical \\
Material & Bulk niobium, RRR $> 300$ \\
Fundamental frequency (TM$_{010}$) & $1.3\;\si{GHz}$ \\
Second mode (TE$_{011}$) & $1.82\;\si{GHz}$ \\
Active length & $1038\;\si{mm}$ \\
Iris diameter & $70\;\si{mm}$ \\
Equator diameter & $206\;\si{mm}$ \\
Wall thickness & $2.8\;\si{mm}$ \\
$Q_0$ (TM$_{010}$, $2\;\si{K}$) & $> 2 \times 10^{10}$ \\
$Q_0$ (TE$_{011}$, $2\;\si{K}$) & $> 1 \times 10^{10}$ \\
Accelerating gradient (TM) & $25\;\si{MV/m}$ \\
Peak surface magnetic field & $105\;\si{mT}$ \\
Peak surface electric field & $50\;\si{MV/m}$ \\
Stored energy per cell (TM) & $255\;\si{J}$ at $25\;\si{MV/m}$ \\
Total stored energy (9-cell TM) & $2295\;\si{J}$ \\
$R/Q$ (TM$_{010}$) & $1036\;\si{\Omega}$ \\
Geometry factor $G$ & $271\;\si{\Omega}$ \\
Operating temperature & $2.0\;\si{K}$ (superfluid He) \\
Surface treatment & Nitrogen-infusion + $120\;^\circ$C bake \\
\bottomrule
\end{longtable}

\subsubsection{Dual-Mode Operation}

Each cavity is simultaneously driven in TM$_{010}$ and TE$_{011}$
modes to break geometry transparency:

\paragraph{RF feed configuration:}
\begin{itemize}[nosep]
\item TM$_{010}$ input coupler: $Q_{\rm ext} = 10^7$ (matched
for $10\;\si{W}$ forward power)
\item TE$_{011}$ input coupler: $Q_{\rm ext} = 5 \times 10^7$
(separate antenna, orthogonal polarization)
\item Field probe (TM): $Q_{\rm probe} = 10^{12}$, for phase
and amplitude monitoring
\item Field probe (TE): $Q_{\rm probe} = 10^{12}$
\item HOM coupler: Standard TESLA design for parasitic mode
damping
\end{itemize}

\paragraph{Phase control:}
\begin{itemize}[nosep]
\item Relative phase $\phi_{\rm TE-TM}$ controlled by digital
IQ modulator
\item Phase resolution: $< 0.001\;\si{rad}$
\item Phase stability: $< 0.01\;\si{rad}$ RMS over $10^4\;\si{s}$
\item Phase modulation at $2\omega$ with programmable depth
$m = 0$--$0.2$
\end{itemize}

\subsubsection{Array Layout}

Four cavities arranged vertically along the chamber axis:

\begin{longtable}{p{0.12\textwidth}p{0.20\textwidth}p{0.25\textwidth}p{0.33\textwidth}}
\toprule
\textbf{Cavity} & \textbf{$z$ position} & \textbf{$\psi$-mode role} & \textbf{Notes} \\
\midrule
\endhead
Cav-1 & $z = 0.0\;\si{m}$ & Antinode 1 & Reference cavity \\
Cav-2 & $z = 0.75\;\si{m}$ & Antinode 2 & Phase-locked to Cav-1 \\
Cav-3 & $z = 1.50\;\si{m}$ & Antinode 3 & Phase-locked to Cav-1 \\
Cav-4 & $z = 2.25\;\si{m}$ & Antinode 4 & Phase-locked to Cav-1 \\
\bottomrule
\end{longtable}

\paragraph{$\psi$-mode tube dimensions:}
\begin{itemize}[nosep]
\item Tube cross-section: $A_\psi = \pi(0.5)^2 = 0.785\;\si{m^2}$
\item Tube height: $L = 3.0\;\si{m}$
\item Total cavity aperture: $A_{\rm cav,tot} = 4 \times \pi(0.035)^2
= 1.54 \times 10^{-2}\;\si{m^2}$
\item Area ratio: $A_{\rm cav,tot}/A_\psi = 0.0196$
\end{itemize}

\subsection{Cryogenic System}

\subsubsection{Helium Plant}

\begin{longtable}{p{0.40\textwidth}p{0.50\textwidth}}
\toprule
\textbf{Parameter} & \textbf{Specification} \\
\midrule
\endhead
Refrigeration capacity at $2\;\si{K}$ & $100\;\si{W}$ \\
Refrigeration capacity at $4.5\;\si{K}$ & $500\;\si{W}$ \\
LHe storage & $2000\;\si{L}$ dewar \\
Compressor power & $150\;\si{kW}$ \\
Cooldown time (300 K to 2 K) & $48\;\si{h}$ \\
Steady-state He consumption & $< 5\;\si{L/h}$ (with recovery) \\
Cryostat vacuum & $< 10^{-4}\;\si{Pa}$ \\
Thermal shield temperature & $80\;\si{K}$ (LN$_2$ or cryo-cooler) \\
\bottomrule
\end{longtable}

\subsubsection{Cryostat Design}

Each SRF cavity sits in an individual cryostat module:
\begin{itemize}[nosep]
\item Inner vessel (He bath): $\varnothing 280\;\si{mm}$ Ti alloy
\item Thermal shield: $80\;\si{K}$ Al with MLI blanket
\item Vacuum jacket: $\varnothing 600\;\si{mm}$ SS
\item Magnetic shield (within cryostat): $\varnothing 260\;\si{mm}$
Nb-Ti or Cryoperm
\item Support: Low-thermal-conductivity G-10 struts
\item Total heat load per module: $\sim 20\;\si{W}$ at $2\;\si{K}$
\end{itemize}

\subsection{RF Power System}

\begin{longtable}{p{0.40\textwidth}p{0.50\textwidth}}
\toprule
\textbf{Parameter} & \textbf{Specification} \\
\midrule
\endhead
TM$_{010}$ RF source & Solid-state amplifier, $1.3\;\si{GHz}$ \\
TM$_{010}$ CW power & $50\;\si{W}$ per cavity (CW) \\
TE$_{011}$ RF source & Solid-state amplifier, $1.82\;\si{GHz}$ \\
TE$_{011}$ CW power & $20\;\si{W}$ per cavity (CW) \\
Phase noise (at $10\;\si{kHz}$ offset) & $< -140\;\si{dBc/Hz}$ \\
Amplitude stability & $< 10^{-6}$ RMS \\
AM modulation bandwidth & DC to $100\;\si{MHz}$ \\
Frequency tuning range & $\pm 500\;\si{kHz}$ (piezo tuner) \\
Pulsed power (Phase 3) & $500\;\si{kW}$ klystron, $10\;\si{ms}$ pulse \\
\bottomrule
\end{longtable}

\subsection{Rotating Mass Source (Calibration)}

\begin{longtable}{p{0.40\textwidth}p{0.50\textwidth}}
\toprule
\textbf{Parameter} & \textbf{Specification} \\
\midrule
\endhead
Material & Tungsten alloy (W-Ni-Cu, $\rho = 18{,}500\;\si{kg/m^3}$) \\
Geometry & Two semi-cylindrical masses on rotor \\
Total rotating mass & $2 \times 204 = 408\;\si{kg}$ \\
Rotor diameter & $300\;\si{mm}$ \\
Rotor length & $300\;\si{mm}$ \\
Rotation frequency & $0$--$60\;\si{Hz}$ (variable speed drive) \\
Balance grade & ISO G0.4 \\
Bearing type & Air bearings (low vibration) \\
Position (relative to sensors) & External to main vessel, $r = 0.8\;\si{m}$
from axis \\
$\psi$ signal at sensor & $\delta\psi \approx 2GM/(c^2 r)
\times \sin(2\pi f_{\rm rot}\,t) \sim 6 \times 10^{-25}\;\si{} \times f_2$
component \\
Gravitational signal & $\delta g \approx 2GM r_0/(r^3)
\sim 8.5 \times 10^{-8}\;\si{m/s^2}$ (static, at $r=0.8\;\si{m}$) \\
\bottomrule
\end{longtable}

%=============================================================================
\section{Sensor Suite Design}
%=============================================================================

\subsection{Sensor Layout}

Sensors are arranged at three stations within the vacuum chamber:

\begin{longtable}{p{0.12\textwidth}p{0.15\textwidth}p{0.63\textwidth}}
\toprule
\textbf{Station} & \textbf{$z$ position} & \textbf{Instruments} \\
\midrule
\endhead
Lower & $z = 0.3\;\si{m}$ & FP cavity \#1, Sr clock MOT,
superconducting gravimeter (vertical),
MZI lower beam \\
Center & $z = 1.5\;\si{m}$ & Superconducting gravimeters
($x,y$ horizontal), atom interferometer zone,
diagnostic optics \\
Upper & $z = 2.7\;\si{m}$ & FP cavity \#2, diagnostic clock
(Yb, Phase 2+), MZI upper beam \\
\bottomrule
\end{longtable}

\subsection{Optical Interferometer (MZI)}

\subsubsection{Detailed Design}

\begin{longtable}{p{0.40\textwidth}p{0.50\textwidth}}
\toprule
\textbf{Parameter} & \textbf{Specification} \\
\midrule
\endhead
Topology & Mach-Zehnder, vertical arm separation \\
Arm length & $1.0\;\si{m}$ (folded if needed) \\
Vertical separation $\Delta h$ & $2.4\;\si{m}$ (between stations) \\
Light source & Nd:YAG $1064\;\si{nm}$, $1\;\si{W}$,
linewidth $< 1\;\si{kHz}$ \\
Beam splitters & Low-loss dielectric, 50:50 \\
Mirrors & Superpolished, $R > 99.99\%$ \\
Readout & Balanced heterodyne at $80\;\si{MHz}$ \\
Phase resolution & $10^{-10}\;\si{rad/\sqrt{Hz}}$ (shot-noise limited) \\
Path length stability & Active fiber stabilization,
$\delta L < 10^{-13}\;\si{m/\sqrt{Hz}}$ \\
Common-mode rejection & $> 10^4$ for laser frequency noise \\
Beam diameter & $3\;\si{mm}$ (to minimize diffraction) \\
Alignment & Automatic beam steering with quad PDs \\
\bottomrule
\end{longtable}

\paragraph{$\psi$ response:}
Differential phase for a uniform $\Delta\psi$ across the
vertical separation:
\begin{equation}
\Delta\phi_{\rm MZI} = \frac{2\pi L}{\lambda}\,\Delta\psi
= \frac{2\pi \times 1.0}{1.064 \times 10^{-6}}\,\Delta\psi
= 5.91 \times 10^6\,\Delta\psi \;\si{rad}.
\end{equation}
Minimum detectable: $\Delta\psi_{\rm min} = 10^{-10}/(5.91 \times 10^6)
\times (1/\sqrt{\tau}) = 1.7 \times 10^{-17}/\sqrt{\tau/\si{s}}$.

\subsection{Fabry-P\'erot Cavity Sensors}

\subsubsection{Cavity Design}

Two identical FP cavities at upper and lower stations:

\begin{longtable}{p{0.40\textwidth}p{0.50\textwidth}}
\toprule
\textbf{Parameter} & \textbf{Specification} \\
\midrule
\endhead
Spacer material & ULE glass (CTE zero-crossing at $\sim 20\;^\circ$C) \\
Length & $100\;\si{mm}$ \\
Mirror substrate & Fused silica, superpolished \\
Mirror coating & Ion-beam-sputtered Ta$_2$O$_5$/SiO$_2$, 40 layers \\
Reflectivity & $R > 99.9985\%$ \\
Finesse & $\mathcal{F} \approx 400{,}000$ \\
Free spectral range & $1.5\;\si{GHz}$ \\
Linewidth & $3.75\;\si{kHz}$ \\
Thermal noise (Brownian) & $\sim 3 \times 10^{-16}$ at $1\;\si{s}$ \\
Fractional frequency stability & $< 10^{-16}$ at $0.1$--$100\;\si{s}$ \\
Laser locking & Pound-Drever-Hall, bandwidth $100\;\si{kHz}$ \\
Laser & ECDL at $698\;\si{nm}$ (shared with Sr clock) \\
Operating temperature & $20.00 \pm 0.001\;^\circ$C \\
Orientation & Vertical (sensitive to $\partial\psi/\partial z$
through differential expansion) \\
\bottomrule
\end{longtable}

\paragraph{$\psi$ response:}
$\delta f/f = -2\,\delta\psi$.
Minimum detectable (at $1\;\si{s}$): $\delta\psi_{\rm min}
= (1/2) \times 10^{-16} = 5 \times 10^{-17}$.

\subsection{Optical Lattice Clock (Sr)}

\begin{longtable}{p{0.40\textwidth}p{0.50\textwidth}}
\toprule
\textbf{Parameter} & \textbf{Specification} \\
\midrule
\endhead
Species & $^{87}$Sr \\
Transition & $^1$S$_0 \to$ $^3$P$_0$ at $698.4\;\si{nm}$
($429.228\;\si{THz}$) \\
Lattice wavelength & $813.4274\;\si{nm}$ (magic wavelength) \\
Lattice depth & $80\,E_r$ ($E_r$ = recoil energy) \\
Atom number & $\sim 10^4$ per cycle \\
Interrogation time & $1$--$10\;\si{s}$ \\
Dick effect limit & $3 \times 10^{-17}/\sqrt{\tau}$ \\
Systematic uncertainty & $< 2 \times 10^{-18}$ (target) \\
BBR shift uncertainty & $< 5 \times 10^{-19}$ (in-vacuum) \\
Clock laser & PDH-locked to FP cavity \#1 \\
Cycle time & $\sim 1.5\;\si{s}$ \\
Location & Lower sensor station \\
\bottomrule
\end{longtable}

\subsection{Superconducting Gravimeters}

Three superconducting gravimeters provide acceleration readout:

\begin{longtable}{p{0.40\textwidth}p{0.50\textwidth}}
\toprule
\textbf{Parameter} & \textbf{Specification} \\
\midrule
\endhead
Type & Nb levitated sphere in persistent-current trap \\
Sensitivity & $< 10^{-12}\;\si{m/s^2/\sqrt{Hz}}$ \\
Bandwidth & $10^{-4}$--$1\;\si{Hz}$ \\
Dynamic range & $\pm 10^{-4}\;\si{m/s^2}$ \\
Long-term drift & $< 10^{-10}\;\si{m/s^2/yr}$ \\
Operating temperature & $4.2\;\si{K}$ \\
Orientations & 1 vertical (center station), 2 horizontal
($x$, $y$, center station) \\
Readout & SQUID, lock-in at pump frequency \\
\bottomrule
\end{longtable}

\subsection{Atom Interferometer (Phase 2+)}

\begin{longtable}{p{0.40\textwidth}p{0.50\textwidth}}
\toprule
\textbf{Parameter} & \textbf{Specification} \\
\midrule
\endhead
Species & $^{87}$Rb \\
Configuration & Mach-Zehnder, vertical \\
Beam-splitting & Raman transitions at $780\;\si{nm}$ \\
Interrogation time $T$ & $0.3\;\si{s}$ \\
Interferometer length & $2T^2 g/2 \approx 0.44\;\si{m}$ \\
Phase sensitivity & $\sim k_{\rm eff}\,a\,T^2
= 1.6 \times 10^7 \times a \times 0.09\;\si{rad}$ \\
Acceleration sensitivity & $\sim 10^{-9}\;\si{m/s^2/\sqrt{Hz}}$ \\
Repetition rate & $\sim 1\;\si{Hz}$ \\
Atom source & 3D MOT + molasses, $10^6$ atoms per shot \\
Location & Center station, vertical path \\
\bottomrule
\end{longtable}

\subsection{Timing Distribution System}

\begin{longtable}{p{0.40\textwidth}p{0.50\textwidth}}
\toprule
\textbf{Parameter} & \textbf{Specification} \\
\midrule
\endhead
Primary reference & Active hydrogen maser,
$\sigma_y \sim 10^{-13}$ at $1\;\si{s}$ \\
Optical distribution & Frequency comb (Er:fiber, $250\;\si{MHz}$
rep rate) \\
Comb stability & Phase-locked to Sr clock laser \\
Fiber distribution & PM fiber, actively stabilized \\
Transfer stability & $< 10^{-19}$ ($10^4\;\si{s}$ integration) \\
RF distribution & $10\;\si{MHz}$, $100\;\si{MHz}$, $1\;\si{GHz}$
synthesized from comb \\
GPS receiver & Multi-constellation, for external comparison \\
PPS output & $< 1\;\si{ns}$ accuracy to UTC \\
\bottomrule
\end{longtable}

%=============================================================================
\section{Control System Architecture}
%=============================================================================

\subsection{Hardware Overview}

\begin{longtable}{p{0.25\textwidth}p{0.65\textwidth}}
\toprule
\textbf{Component} & \textbf{Specification} \\
\midrule
\endhead
FPGA platform & Xilinx Ultrascale+ (or equivalent), $500\;\si{MHz}$
clock \\
ADC (fast) & 16-channel, 16-bit, $250\;\si{MSPS}$ \\
ADC (precision) & 32-channel, 24-bit, $1\;\si{kSPS}$ \\
DAC (fast) & 8-channel, 16-bit, $1\;\si{GSPS}$ \\
DAC (slow) & 16-channel, 20-bit, $100\;\si{kSPS}$ \\
Digital I/O & 64 channels \\
Timing & White Rabbit switches ($< 1\;\si{ns}$ sync) \\
Computing & Dual-CPU server with GPU (for offline Kalman filter) \\
Network & $10\;\si{Gbit}$ Ethernet backbone \\
Storage & RAID-6 array, $100\;\si{TB}$ \\
\bottomrule
\end{longtable}

\subsection{Real-Time State Estimator}

\subsubsection{Kalman Filter Design}

\paragraph{State vector:}
The $\psi$ field is discretized on a cylindrical grid
($N_r = 10$, $N_z = 30$, $N_\theta = 1$ for azimuthal symmetry),
giving $N = 300$ spatial nodes.  The state vector is:
\begin{equation}
\mathbf{x} = [\psi_1, \ldots, \psi_N,
\dot{\psi}_1, \ldots, \dot{\psi}_N]^T
\in \mathbb{R}^{600}.
\end{equation}

\paragraph{Process model:}
The linearized $\psi$ field equation in matrix form:
\begin{equation}
\dot{\mathbf{x}} = \mathbf{A}\,\mathbf{x} + \mathbf{B}\,\mathbf{u}
+ \mathbf{w},
\end{equation}
where $\mathbf{A}$ encodes the discretized field equation
(including screening), $\mathbf{B}\,\mathbf{u}$ is the EM
source driving, and $\mathbf{w}$ is process noise.

\paragraph{Measurement model:}
\begin{equation}
\mathbf{z} = \mathbf{H}\,\mathbf{x} + \mathbf{v},
\end{equation}
where $\mathbf{H}$ maps the field state to:
\begin{itemize}[nosep]
\item MZI phase: $z_{\rm MZI} = (2\pi L/\lambda)\sum_i w_i\,\psi_i$
\item FP frequency: $z_{\rm FP,k} = -2\,\psi(z_k)$
\item Clock ratio: $z_{\rm clock} = (K_A^{\rm obs} - K_B^{\rm obs})
\,\psi(z_{\rm clock})$
\item Gravimeters: $z_{{\rm grav},k} = (c^2/2)\,(\partial\psi/\partial z)
\big|_{z_k}$
\end{itemize}

\paragraph{Update rate:} $10\;\si{kHz}$ (FPGA-implemented
for fast channels), $1\;\si{Hz}$ (CPU for clock channels).

\subsection{Feedback Controller}

\subsubsection{Control Architecture}

Two nested feedback loops:

\paragraph{Inner loop (fast, FPGA):}
\begin{itemize}[nosep]
\item Controls: TE/TM relative phase $\phi$, RF amplitude,
cavity frequency (via piezo tuner)
\item Bandwidth: $1\;\si{kHz}$
\item Objective: Maintain EM drive at $2\omega = \Omega_\psi$
resonance
\end{itemize}

\paragraph{Outer loop (slow, CPU):}
\begin{itemize}[nosep]
\item Controls: Modulation depth $m$, frequency scan rate,
integration time
\item Bandwidth: $0.01\;\si{Hz}$
\item Objective: Optimize $\psi$ generation or scan frequency
parameter space
\end{itemize}

\subsubsection{Operating Modes}

\begin{longtable}{p{0.18\textwidth}p{0.72\textwidth}}
\toprule
\textbf{Mode} & \textbf{Description} \\
\midrule
\endhead
\texttt{SCAN} & Swept-frequency search for $\psi$ modes.  Drive
frequency steps through programmed range.  Lock-in detection at
each step. \\
\texttt{LOCK} & Frequency locked to identified $\psi$ mode.  PLL
maintains $2\omega = \Omega_\psi$.  Long integration for
sensitivity. \\
\texttt{SHAPE} & Active $\psi$-field shaping.  Multiple cavities
driven with phase offsets to produce desired spatial profile. \\
\texttt{CALIBRATE} & Rotating mass or known gravitational source.
Validates sensor chain end-to-end. \\
\texttt{NULL} & All sources off.  Measures ambient background
for subtraction. \\
\texttt{SAFE} & All RF off.  Monitoring only.  Default power-on
state. \\
\bottomrule
\end{longtable}

%=============================================================================
\section{Safety System Design}
%=============================================================================

\subsection{Safety Philosophy}

The safety system is designed to the IEC 61508 SIL-2 standard
(Safety Integrity Level 2), providing:
\begin{itemize}[nosep]
\item Independent safety-rated sensors (separate from measurement
sensors)
\item Hardwired interlock logic (not software-dependent)
\item Fail-safe design (loss of power $\to$ safe state)
\item Dual-redundant trip channels
\end{itemize}

\subsection{Safety Instrumented System (SIS)}

\subsubsection{Safety Sensors}

\begin{longtable}{p{0.20\textwidth}p{0.25\textwidth}p{0.22\textwidth}p{0.23\textwidth}}
\toprule
\textbf{Sensor} & \textbf{Measurand} & \textbf{Range} & \textbf{Response} \\
\midrule
\endhead
Safety accel.~1 & $|a_z|$ & $0$--$1\;\si{m/s^2}$ & $< 1\;\si{ms}$ \\
Safety accel.~2 & $|a_z|$ & $0$--$1\;\si{m/s^2}$ & $< 1\;\si{ms}$ \\
Safety MZI & $|\Delta\phi|$ & $0$--$2\pi$ & $< 10\;\si{\mu s}$ \\
Personnel sensor & Exclusion zone & Presence/absence & $< 100\;\si{ms}$ \\
RF power mon. & Forward power & $0$--$1\;\si{MW}$ & $< 1\;\si{\mu s}$ \\
Cryo sensors & He level, pressure & Multiple & $< 1\;\si{s}$ \\
Vacuum gauge & Chamber pressure & $10^{-8}$--$10^5\;\si{Pa}$ & $< 1\;\si{s}$ \\
O$_2$ monitor & Ambient O$_2$ & $0$--$25\%$ & $< 5\;\si{s}$ \\
\bottomrule
\end{longtable}

\subsubsection{Interlock Logic}

\begin{tcolorbox}[colback=red!3, colframe=red!60!black,
title=Safety Interlock Trip Matrix]
\begin{verbatim}
CONDITION                    ACTION
---------                    ------
|a_anom| > 10^-4 m/s^2  --> RF KILL + CAVITY DETUNE
|a_anom| > 10^-6 m/s^2  --> WARNING + LOG
|dpsi/dt| > 10^-6 /s    --> RF KILL + CAVITY DETUNE
Vacuum > 10^-3 Pa       --> RF OFF (protect cavities)
He level < 50%           --> RF POWER REDUCE
He level < 20%           --> RF OFF
Personnel in zone        --> RF INTERLOCK
O2 < 19.5%              --> ALARM + VENTILATION
Any 2 safety sensors     --> FULL SHUTDOWN (Phase 4)
  disagreeing
\end{verbatim}
\end{tcolorbox}

\subsubsection{Shutdown Sequence}

\begin{longtable}{p{0.10\textwidth}p{0.18\textwidth}p{0.62\textwidth}}
\toprule
\textbf{Phase} & \textbf{Time} & \textbf{Action} \\
\midrule
\endhead
1 & $< 1\;\si{\mu s}$ & PIN diode switches kill RF drive to all cavities \\
2 & $< 100\;\si{\mu s}$ & Fast piezo tuners detune cavities by $> 100$
bandwidths \\
3 & $< 10\;\si{ms}$ & Stored energy rings down ($\tau = Q/(\pi f)
\sim 2.4\;\si{s}$; detuning accelerates) \\
4 & $< 1\;\si{s}$ & Mains power disconnect via safety-rated
contactors \\
5 & $< 60\;\si{s}$ & All systems to monitored safe state \\
\bottomrule
\end{longtable}

\subsection{Personnel Safety}

\begin{itemize}[nosep]
\item $5\;\si{m}$ exclusion zone during high-power operation,
enforced by light curtains and interlocked doors
\item Personal dosimeters: wearable MEMS accelerometers logging
at $100\;\si{Hz}$, alarm at $10^{-4}\;\si{m/s^2}$
\item RF exposure: maintained below ICNIRP occupational limits
at all accessible locations
\item Cryogenic: ODH (Oxygen Deficiency Hazard) monitoring,
forced ventilation, emergency breathing apparatus
\item Magnetic: Access restricted during degaussing operations
\item Radiation: No significant ionizing radiation sources;
standard lab monitoring
\end{itemize}

%=============================================================================
\section{Expected Performance and Sensitivity}
%=============================================================================

\subsection{Design Law Evaluation}

Using the design law (Eq.~\ref{eq:design-law-ref}):
\begin{equation}
|\lambda - 1|_{\rm min} = \frac{\pi\gamma_\psi}{c_s\,U_0\,m}
\frac{A_\psi^2}{\kappa_{\rm eff}\,A_{\rm cav,tot}}
\label{eq:design-law-ref}
\end{equation}

\begin{longtable}{p{0.25\textwidth}p{0.18\textwidth}p{0.18\textwidth}p{0.18\textwidth}}
\toprule
\textbf{Parameter} & \textbf{Phase 1} & \textbf{Phase 2} & \textbf{Phase 3} \\
\midrule
\endhead
$U_0$ (J) & $9 \times 10^3$ & $9 \times 10^3$ & $10^6$ \\
$m$ & $0.01$ & $0.05$ & $0.10$ \\
$A_\psi$ (m$^2$) & $0.785$ & $0.785$ & $0.27$ \\
$A_{\rm cav,tot}$ (m$^2$) & $0.0154$ & $0.0154$ & $0.030$ \\
$\gamma_\psi/\Omega_\psi$ & $10^{-3}$ & $10^{-3}$ & $10^{-3}$ \\
$\kappa_{\rm eff}$ & $1$ & $1$ & $1$ \\
$c_s$ & $c$ & $c$ & $c$ \\
\midrule
$|\lambda{-}1|_{\rm min}$ & $2.7 \times 10^{-7}$ & $5.3 \times 10^{-8}$
& $5.1 \times 10^{-12}$ \\
\bottomrule
\end{longtable}

\paragraph{Note:} These are the parametric instability thresholds.
The driven-resonance channel with TE+TM superposition provides
an independent sensitivity estimate (see main paper
Eq.~\ref{eq:psi-amplitude}).  With lock-in detection and long
integration, the effective reach extends further as shown in
the main paper Table~\ref{tab:lambda-reach}.

\subsection{Noise Budget Summary}

\begin{longtable}{p{0.30\textwidth}p{0.20\textwidth}p{0.20\textwidth}p{0.20\textwidth}}
\toprule
\textbf{Noise Source} & \textbf{Level} & \textbf{$\psi$ equiv.} & \textbf{Regime} \\
\midrule
\endhead
Shot noise (MZI) & $3 \times 10^{-10}\;\si{rad/\sqrt{Hz}}$
& $5 \times 10^{-17}/\sqrt{\si{Hz}}$ & White \\
Seismic (after isolation) & $10^{-9}\;\si{m/s^2/\sqrt{Hz}}$
& $2 \times 10^{-26}/\sqrt{\si{Hz}}$ & $1/f$ \\
Thermal (FP cavity) & $3 \times 10^{-16}$
& $1.5 \times 10^{-16}$ & $1/\sqrt{f}$ \\
Laser frequency & $10^{-14}/\sqrt{\si{Hz}}$
& CMR $> 10^4$ & White \\
Residual gas & $10^{-12}\;\si{rad/\sqrt{Hz}}$
& $2 \times 10^{-19}/\sqrt{\si{Hz}}$ & White \\
Clock QPN & $5 \times 10^{-17}/\sqrt{\tau}$
& species-dependent & White \\
Gravimeter thermal & $10^{-12}\;\si{m/s^2/\sqrt{Hz}}$
& $2 \times 10^{-29}/\sqrt{\si{Hz}}$ & White \\
ADC quantization & $< 10^{-7}$ (24-bit)
& negligible & White \\
\midrule
\textbf{Combined floor} & --- & $\sim 5 \times 10^{-17}/\sqrt{\si{Hz}}$
& Shot-noise limited \\
\bottomrule
\end{longtable}

\subsection{Sensitivity Timeline}

\begin{longtable}{p{0.20\textwidth}p{0.18\textwidth}p{0.22\textwidth}p{0.30\textwidth}}
\toprule
\textbf{Phase} & \textbf{Duration} & \textbf{$\Delta\psi_{\rm min}$}
& \textbf{$|\lambda{-}1|_{\rm min}$} \\
\midrule
\endhead
Phase 1, scan & 12 days & $10^{-15}$ per step & $10^{-8}$ (broadband) \\
Phase 1, lock & 1 month & $5 \times 10^{-17}$ & $10^{-9}$ \\
Phase 2, lock & 3 months & $10^{-17}$ & $10^{-11}$ \\
Phase 3, lock & 6 months & $10^{-19}$ & $10^{-14}$ \\
\bottomrule
\end{longtable}

%=============================================================================
\section{Commissioning and Operations}
%=============================================================================

\subsection{Commissioning Sequence}

\begin{enumerate}[nosep]
\item \textbf{Mechanical installation} (4 weeks): Vacuum vessel,
isolation platform, shielding, cryostats.
\item \textbf{Vacuum qualification} (2 weeks): Leak test,
pumpdown, bakeout, RGA survey.
\item \textbf{Cryogenic commissioning} (2 weeks): Cooldown,
thermal stability verification, He plant tuning.
\item \textbf{RF commissioning} (2 weeks): Individual cavity
conditioning, $Q_0$ measurement, dual-mode excitation test.
\item \textbf{Sensor commissioning} (4 weeks): MZI alignment,
FP cavity locking, gravimeter cooldown, clock first light.
\item \textbf{Integration test} (2 weeks): Full system operation,
data pipeline validation, rotating mass calibration.
\item \textbf{Safety validation} (1 week): Trip test matrix,
shutdown timing verification, personnel safety audit.
\end{enumerate}

Total commissioning: $\sim 17$ weeks.

\subsection{Operational Procedures}

\subsubsection{Daily Operations}

\begin{enumerate}[nosep]
\item Verify vacuum, cryogenic, and environmental parameters
within specification.
\item Check all sensor calibrations against internal references.
\item Execute scheduled measurement program (\texttt{SCAN},
\texttt{LOCK}, or \texttt{CALIBRATE}).
\item Monitor safety system status continuously.
\item Perform end-of-day data quality review.
\end{enumerate}

\subsubsection{Weekly Maintenance}

\begin{enumerate}[nosep]
\item LHe dewar refill or recovery system check.
\item Mu-metal degaussing cycle.
\item Sensor recalibration against rotating mass source.
\item Vibration isolation performance verification.
\item Data backup and transfer to analysis cluster.
\end{enumerate}

%=============================================================================
\section{Cost and Schedule Estimates}
%=============================================================================

\subsection{Phase 1 Capital Cost}

\begin{longtable}{p{0.45\textwidth}p{0.20\textwidth}p{0.25\textwidth}}
\toprule
\textbf{Item} & \textbf{Cost (USD)} & \textbf{Notes} \\
\midrule
\endhead
Vacuum vessel + shielding & \$350K & Custom fabrication \\
Vacuum pumps + gauges & \$120K & Commercial \\
Vibration isolation & \$200K & Commercial + custom \\
SRF cavities (4$\times$) & \$800K & Niobium + fabrication \\
Cryostats (4$\times$) & \$400K & Custom \\
He refrigerator & \$600K & Commercial \\
RF power systems & \$300K & Solid-state amplifiers \\
Optical interferometer & \$150K & Custom optics \\
FP cavities (2$\times$) & \$200K & Custom \\
Sr clock system & \$500K & Laser + vacuum + optics \\
Superconducting gravimeters (3$\times$) & \$600K & Commercial \\
Control electronics + DAQ & \$250K & FPGA + ADC + DAC \\
Timing (maser + comb) & \$350K & Commercial \\
Rotating mass source & \$80K & Custom machining \\
Installation + commissioning & \$400K & Labor + consumables \\
\midrule
\textbf{Phase 1 Total} & \textbf{\$5.3M} & \\
\bottomrule
\end{longtable}

\subsection{Operating Cost (Annual)}

\begin{longtable}{p{0.45\textwidth}p{0.20\textwidth}}
\toprule
\textbf{Item} & \textbf{Cost (USD/yr)} \\
\midrule
\endhead
Liquid helium & \$150K \\
Electricity ($200\;\si{kW}$ avg) & \$175K \\
Personnel (3 FTE) & \$450K \\
Consumables + maintenance & \$100K \\
\midrule
\textbf{Annual Total} & \textbf{\$875K} \\
\bottomrule
\end{longtable}

\subsection{Schedule}

\begin{longtable}{p{0.12\textwidth}p{0.20\textwidth}p{0.58\textwidth}}
\toprule
\textbf{Month} & \textbf{Milestone} & \textbf{Deliverable} \\
\midrule
\endhead
M0--M3 & Detailed design & Final drawings, procurement specifications \\
M3--M9 & Procurement & Long-lead items (SRF cavities, He plant) \\
M6--M12 & Installation & Vacuum vessel, isolation, cryogenics \\
M12--M16 & Commissioning & Per Sec.~10.1 sequence \\
M16--M18 & Phase 1 science & Broadband $\psi$-mode scan \\
M18--M24 & Phase 1 results & First $|\lambda - 1|$ constraint \\
M24--M36 & Phase 2 upgrade & Clock + atom interferometer integration \\
M36+ & Phase 3 planning & Pulsed high-power design \\
\bottomrule
\end{longtable}

%=============================================================================
\section{Conclusion}
%=============================================================================

This engineering design document provides a complete,
deployment-ready specification for the Laboratory Scalar Density
Field Generator.  Key features:

\begin{enumerate}[nosep]
\item \textbf{Proven technology base}: SRF cavities, optical
clocks, and superconducting gravimeters are mature technologies.
\item \textbf{Multi-modal detection}: Independent optical, atomic,
and inertial measurement channels provide cross-validation.
\item \textbf{Systematic uncertainty control}: Triple-layer
magnetic shielding, three-stage thermal stabilization,
multi-stage vibration isolation.
\item \textbf{Safety by design}: Hardware-enforced limits with
microsecond response times.
\item \textbf{Phased approach}: Phase 1 achieves 4 orders of
magnitude improvement over current bounds with existing technology.
\item \textbf{Cost-effective}: Phase 1 capital cost of \$5.3M is
comparable to a single precision metrology experiment.
\end{enumerate}

The apparatus is the enabling technology for all DFD patent
applications and represents the critical first step toward
laboratory validation of scalar density field dynamics.

\end{document}
